\section{Vérification}
\label{sec:verification}

\subsection{Méthodes de vérification}

\begin{table}[H]
  \centering\small
  \begin{tabularx}{\linewidth}{@{}lp{28mm}X@{}}
    \toprule
    \textbf{Code} & \textbf{Méthode} & \textbf{Définition} \\
    \midrule
    I & Inspection    & Examen de l'application, de ses écrans ou de sa documentation, sans la faire fonctionner. \\
    A & Analyse       & Démonstration par le calcul, la modélisation ou la revue de conception. \\
    D & Démonstration & Fonctionnement de l'application observé, sans instrumentation ni mesure chiffrée. \\
    T & Test          & Fonctionnement de l'application dans des conditions contrôlées, résultat comparé à un critère. \\
    \bottomrule
  \end{tabularx}
  \caption{Méthodes de vérification}
\end{table}

\subsection{Niveaux et moyens de vérification}

\begin{table}[H]
  \centering\small
  \begin{tabularx}{\linewidth}{@{}p{32mm}p{52mm}X@{}}
    \toprule
    \textbf{Niveau} & \textbf{Moyens} & \textbf{Portée} \\
    \midrule
    Émulateur Pebble & Application lancée sur l'émulateur de chacune des six plateformes, avec une maison simulée (annexe, §\ref{sec:maison-essai}). & Écrans, navigation, lisibilité, adaptation à chaque modèle (couleur, noir et blanc, rond). \\
    Compte Tado d'essai & Application reliée à un vrai compte Tado, sur une maison dont on connaît l'état. & Lecture de la situation, commandes, présence, messages d'erreur, liaison et rupture de liaison. \\
    Recette à la maison & Une montre réelle portée une semaine dans la maison de l'utilisateur. & Usage au poignet, nombre d'appuis, fraîcheur des données. \\
    \bottomrule
  \end{tabularx}
  \caption{Niveaux de vérification}
\end{table}

La maison d'essai décrite en annexe (§\ref{sec:maison-essai}) sert de jeu d'essai commun~:
elle contient une pièce qui chauffe fort, une pièce forcée, une pièce sans chauffe récente,
une pièce froide et une humidité mesurée.

\acompleter{niveau retenu pour chaque exigence, critères d'acceptation chiffrés
(nombre d'appuis, âge maximal des données, fréquence d'actualisation), référence du test qui
la vérifie.}

\subsection{Matrice de vérification}

La matrice ci-dessous est générée à partir des exigences du §3.

\matriceverification
